\section{Discussion}
\label{sec:discussion}

\subsection{Answers to the research questions}

By definition of the frozen routing relation, RQ1 permits routing exactly when
the eligible source-profile set in \cref{eq:eligible-routes} has cardinality
one. Missing or conflicting evidence is not overridden by precedence, and
routing success does not authorize a candidate. This is a design-level
analytical condition, not an empirical effect.

RQ2 receives a conditional analytical answer within the formal model.
Deliverability requires the conjunction in \cref{eq:extended-acceptance},
followed by the rule that only a clean result may carry bytes and by the
non-overwriting publication transition. The
reconstruction order prevents a structural candidate from satisfying a policy
that requires semantic reconstruction. Under parser soundness and the stated
filesystem assumptions, the conjunction is sufficient for canonical-language
membership and non-publication of a new target for every non-clean outcome.
Whether the implementation preserves each conjunct is an RQ3 question and is
not established by the analytical result.

RQ3 remains unanswered at its full declared scope. The initial bounded
experiment in \cref{sec:server-experiment} covered 50 S0, 30 S1, and 24 inert
S6 generated records. A later controlled-readiness cycle preserved that
104-record baseline and added a 25-record generated scaffold across S2--S5 and
S7. The follow-up recorded 80/80 fidelity and differential checks, 312/312
paired-ablation rows with no pipeline error, and 14/14 fault cases. Removing the
input prescan or supplied-type evidence caused no decision transition in the
baseline rehearsal; lowering the configurable reconstruction floor changed
8/48 applicable accepts to blocks. These are descriptive generated-corpus
observations, not causal mechanism estimates. The full protocol-budget manifest
was not frozen, independent signatures and admitted sources remained zero,
provenance failed on a shared host, and neither the 21 fuzz target-hours nor the
10 controlled performance sessions ran. The complete gate therefore remained
false.

\subsection{Scientific novelty and relation to prior work}

Canonical and complete parsing, language-oriented security, content disarm and
reconstruction, parser differential analysis, supervised media conversion, and
non-overwriting file operations all have established antecedents
\cite{samuel2012,jana2012,nail2014,everparse2019,darpaSafeDocs,belkind2023,
dubin2023,dubin2024ole,koch2025}. Against the selected lines compared in
\cref{tab:research-position}, the method studied here combines the eight
obligations F1--F8 in one authorization relation. Source evidence must select
one route. Policy registers the permitted source-to-output conversion and
minimum reconstruction level before dispatch. An authorizing parser evaluates
the complete candidate on a code path separate from construction. Scanning and
resource checks apply to the emitted object, and only a clean result may reach
the non-overwriting publication operation. This composition makes structural
container closure distinguishable from semantic source-payload independence and
prevents failed reconstruction from becoming pass-through delivery. It does not
claim that the constituent parser, codec, scanner, or filesystem primitives are
individually new.

The resulting security statement is narrower than ``malware-free media.''
Complete parsing excludes unaccounted candidate prefixes and suffixes under the
named grammar. Output evidence convergence makes the generated name, declared
type, byte recognition, and authorizing report agree. On semantic routes, no
source payload is copied or referenced as an output component. Under the stated
filesystem assumptions, a non-clean result returns no deliverable bytes, creates
no new target through the component API, and leaves any pre-existing target
unchanged. These properties neither establish universal harmlessness nor rule
out unknown decoder flaws.

\subsection{Security interpretation}

The distinction between structural and semantic reconstruction determines how
format coverage may be interpreted. As
\cref{tab:format-profile-comparison} records, structural routes may retain
frames, packets, samples, or pixel data while still enforcing container,
metadata, and boundary properties. A threat model that includes codec payloads
must require semantic reconstruction or block the profile. A single
``sanitized'' label would conceal this distinction.

The taxonomy in \cref{sec:background} supports scenario interpretation but does
not determine verdicts. The method measures carrier structure, ambiguity,
metadata, resource behavior, reconstruction, and publication. Threat-campaign
scenarios connect those observations to realistic delivery chains without
inferring malware family or operational role.

\subsection{Deployment implications and validity}

Deployment begins at the local transition where externally influenced media
would otherwise become eligible for storage, decoding, preview, export, or
opening by another application. The upstream system may authenticate a server
or sender, preserve transport integrity, or record provenance, but those
properties do not replace local output authorization. In end-to-end encrypted
messaging, this placement also keeps plaintext and verdicts away from a blind
relay \cite{unger2015,cohngordon2020,abelson2024}. In web, mail, cloud, and
direct-import settings, the same method applies under different provenance
assumptions. Reproducible deployment requires identified parser, policy, and
external-tool versions. Different policies may legitimately reject different
objects.

The primary internal-validity risk is oracle dependence. Expected actions in
the bounded run were generated with the fixtures, so 104/104 agreement measures
fixture-policy consistency rather than independent ground truth. A canonical
parser can share a misunderstanding with a handler. External tools can also
share a codec library. The external checks corroborate emitted format or profile
recognition but do not independently validate expected actions or expected
blocks. Common-mode faults therefore remain possible. The generated corpus can
also overrepresent regular fixtures and underrepresent unusual benign media from
specific devices. Grouping descendants by base artifact prevents direct
mutation leakage across partitions, while held-out device and application
strata are still needed for external validity.

Fidelity creates another trade-off. Stronger semantic reconstruction reduces
source dependence but can change compression artifacts, frame timing, color,
audio characteristics, and container features. The security policy and fidelity
policy must remain separate so that perceptual similarity does not authorize
unsafe output. A profile can be secure under its structural claim and still be
operationally unacceptable, or it can satisfy a perceptual metric while failing
canonical closure. The \(G_{\mathrm{sci}}\) gate rejects both cases for
different reasons.

Availability claims remain local. The formal model and implementation are
designed to bound bytes, dimensions, expansion, elements, deadlines, and
supervised child processes for one invocation. The readiness cycle added four
generated S5 scaffolds and executed 14/14 registered fault cases, but this
under-budget, non-admitted corpus does not establish the bounds over the target
population. The method also does not provide queue-level fairness, tenant
isolation, or network admission control. Memory-safety guarantees narrow the
implementation defect surface but do not prove arithmetic, grammar, or semantic
correctness.

External-codec supervision likewise does not sandbox arbitrary decoder code.
Where decoder compromise is in scope, the semantic routes require the additional
platform isolation stated in \cref{tab:platform-containment}. Without it, the
article makes no endpoint execution-integrity claim against decoder compromise.

\subsection{Evidence state and research directions}

Current evidence comprises implementation traceability, the release-ineligible
104-record server experiment, a 129-record S0--S7 readiness inventory, a
ten-file public-benign compatibility pilot, and a 316-object candidate
inventory. The readiness cycle closed the earlier measurement-level fidelity
failures under unchanged thresholds, completed the baseline paired-ablation and
14-case fault rehearsals, and regenerated six available outputs
byte-identically. A mechanical packet verified all candidate byte identities
and retained 948 inspection outputs. Without an independent registry or
signatures, selection still returned 0/246 sources across 22 cells. The
readiness archive makes the remaining work executable but does not create an
audited corpus record or support efficacy or release claims.

Next, the write-disabled handoff archive must be delivered outside the audit
host to reviewers with distinct key lineages and independence clusters.
Matching signed decisions must precede selection, S0/S1 assignment, full-file
admission, and a protocol-budget S0--S7 manifest. The controlled run must then
move to a single-purpose host, execute the 21-hour fuzz and 10-session
performance schedules with 10,000 hierarchical bootstrap replicates, and
regenerate the complete evidence package. The completed rehearsal fault and
ablation schedules must be rerun against that admitted frozen corpus.
